\section{Limits and falsification}

\subsection{Scope limits}

\textbf{Domain specificity.} The benchmark uses marine science queries. Extension to other domains
requires updating the science marker regex for that domain's technical vocabulary.

\textbf{Model dependence.} Results were produced with gemini-3.5-flash (2026-06-24). A model with
stronger domain-specific knowledge may generate fewer science claims.

\textbf{$\Runcited$ does not measure semantic accuracy.} A claim that is cited to a bound artifact
but semantically wrong is not captured. Extension to factual accuracy requires a reference
ground truth.

\textbf{Step-log quality dependence.} $\Runcited = 0\%$ requires a complete, structured step-log;
sparse or malformed traces would not achieve the same result.

\subsection{Falsification conditions}

\begin{center}
\small
\begin{tabular}{@{}p{4.5cm}p{3.8cm}@{}}
\toprule
Claim & Falsifying condition \\
\midrule
Step-log injection achieves $\Runcited = 0\%$ & Step-log-injected query produces uncovered science claims \\
Hierarchy is stable across query types & Different Scenario~4 query yields high $\Runcited$ \\
Gate context achieves no lift & Gate data injection covers generated science claims \\
$\Runcited$ diagnoses architecture quality & Two architecturally different systems produce identical hierarchies \\
\bottomrule
\end{tabular}
\end{center}

\subsection{Extension path}

Multi-domain validation: run the benchmark on pedestrian routing (pax), clinical care, and
manufacturing. World model evaluation: apply the benchmark to JEPA-based world models
(V-JEPA~\cite{arxiv240408471}) and measure whether planning trace injection reduces $\Runcited$.
Semantic accuracy addition: add a second metric for whether cited claims are factually accurate
relative to their bound artifacts.

\subsection{The benchmark gap in existing physical-world evaluation}

Physical-world AI benchmarks evaluate task completion or prediction accuracy without auditing
intermediate claim binding. PHYRE~\cite{arxiv190805656} measures physical reasoning via puzzle
solve rate. STAR~\cite{arxiv240509711} evaluates situated reasoning in real-world videos via
question answering accuracy. ThreeDWorld~\cite{arxiv210314025} measures embodied planning by
object transport success rate. TravelPlanner~\cite{arxiv240201622} evaluates language agent
planning by plan validity. The comprehensive Embodied AI survey~\cite{arxiv240706886} categorizes
evaluation as navigation, interaction, manipulation, and reasoning accuracy. None measures whether
intermediate reasoning claims are bound to their generating sensor observations. EWOK~\cite{arxiv240509605}
tests world modeling via plausibility matching. The need for evaluation primitives beyond functional
task success has a precedent: human-robot interaction required psychological benchmarks to capture
what task-completion metrics omit~\cite{kahn2007benchmarks}. $\Runcited$ provides the missing
primitive for evidence binding.

\subsection{World model evaluation as a future extension}

A natural extension is to apply this benchmark to world model planning systems whose design
explicitly requires sensor-grounded intermediate representations. LeCun's autonomous machine
intelligence architecture~\cite{arxiv230602572} specifies that the world model module must
learn to represent the information content of the percept that is relevant for predicting
future states --- a grounding requirement on intermediate representations. Joint-embedding
predictive architectures, from I-JEPA~\cite{arxiv230108243} to V-JEPA~\cite{arxiv240408471}, and
Image World Models~\cite{arxiv240300504} are evaluated on downstream accuracy on image and video
benchmarks; no claim-level evidence binding metric is proposed.
Whether $\Runcited$ on JEPA planning traces reproduces Section~5's diagnostic hierarchy is an open
empirical question. Section~4's methodology requires only a structured planning trace; running it on
world model outputs needs no architectural change.
